% Part: computability
% Chapter: recursive-functions
% Section: primes

\documentclass[../../../include/open-logic-section]{subfiles}

\begin{document}

\olfileid[pt]{cmp}{rec}{pri}
\olsection{Números primos}

A quantificação limitada e a minimização limitada nos fornecem um bom
conjunto de recursos para mostrar que funções e relações naturais são
recursivas primitivas. Por exemplo, considere a relação ``$x$ divide
$y$'', escrita $x \mid y$. A relação $x \mid y$ vale quando é possível
dividir $y$ por~$x$ sem resto, isto é, quando $y$ é um múltiplo inteiro
de~$x$. (Quando ela não vale, isto é, quando o resto da divisão de $x$
por $y$ é $> 0$, escrevemos $x \nmid y$.) Em outras palavras, $x \mid y$
se, e somente se, para algum~$z$, $x \cdot z = y$. Evidentemente, tal
$z$, se existir, deve ser $\leq y$. Assim, $x \mid y$ se, e somente se,
para algum $z \le y$, $x \cdot z = y$. Podemos definir a relação
$x \mid y$ por quantificação existencial limitada a partir de $=$ e da
multiplicação:
\[
x \mid y \defiff \bexists{z \leq y}{(x \cdot z) = y}.
\]
Mostramos, portanto, que $x \mid y$ é recursiva primitiva.

Um número natural~$x$ é \emph{primo} se não é nem $0$ nem $1$ e só é
divisível por $1$ e por ele mesmo. Em outras palavras, os números primos
são tais que, sempre que $y \mid x$, ou $y = 1$ ou~$y=x$. Para verificar
se $x$~é primo, basta verificar se $y \mid x$ para todo $y \le x$, pois,
se $y > x$, então automaticamente~$y \nmid x$. Logo, a relação
$\fn{Prime}(x)$, que vale se, e somente se, $x$ é primo, pode ser definida
por
\[
\fn{Prime}(x) \defiff x \geq 2 \land \bforall{y \leq x}{(y \mid x \lif y
  = 1 \lor y = x)}
\]
e, portanto, é recursiva primitiva.

Os números primos são $2$, $3$, $5$, $7$, $11$ etc. Considere a função
$p(x)$ que devolve o $x$-ésimo primo dessa sequência, isto é, $p(0) =
2$, $p(1) = 3$, $p(2) = 5$ etc. (Por conveniência, muitas vezes
escreveremos $p(x)$ como $p_x$ ($p_0=2$, $p_1=3$ etc.)

Se tivéssemos uma função
$\fn{nextPrime(x)}$ que devolvesse o primeiro número primo maior
que~$x$, seria fácil definir~$p$ por recursão primitiva:
\begin{align*}
  p(0) & = 2\\
  p(x+1) & = \fn{nextPrime}(p(x))
\end{align*}
Como $\fn{nextPrime}(x)$ é o menor $y$ tal que $y > x$ e $y$~é primo,
ela pode ser facilmente calculada por busca ilimitada. Mas também pode
ser definida por minimização limitada, graças a um resultado de Euclides:
sempre existe um número primo entre $x$ e $\fact{x}+1$.
\[
  \fn{nextPrime(x)} =
  \bmin{y \leq \fact{x}+1}{(y > x \land \fn{Prime}(y))}.
\]
Isso mostra que $\fn{nextPrime}(x)$ e, portanto, $p(x)$ são recursivas
primitivas (e não apenas computáveis).

(Se você estiver curioso, eis uma demonstração rápida do teorema de
Euclides. Suponha que $p_n$ seja o maior primo $\le x$ e considere o
produto $p = p_0\cdot p_1 \cdot \dots \cdot p_n$ de todos os
primos~$\le x$. Ou $p+1$ é primo, ou existe um primo entre $x$
e~$p+1$. Por quê? Suponha que $p+1$ não seja primo. Então algum número
primo $q \mid p+1$, com $q < p+1$. Nenhum dos primos $\le x$ divide
$p+1$. (Pela definição de~$p$, cada um dos primos $p_i \le x$
divide~$p$, isto é, com resto~$0$. Logo, cada um dos primos
$p_i \le x$ divide $p+1$ com resto~$1$ e, portanto,
$p_i \nmid p+1$.) Assim, $q$ é um primo $> x$ e $< p+1$. E
$p \le \fact{x}$, portanto existe um primo $> x$ e $\le
\fact{x}+1$.)

\begin{prob}
Defina a divisão inteira $d(x, y)$ usando minimização limitada.
\end{prob}

\end{document}
